Au début, chacune des six boîtes $B_1, B_2, B_3, B_4, B_5, B_6$ contient un jeton. Deux types d'opération
sont possibles :
\begin{itemize}
  \item[\emph{Type} 1 :] Choisir une boîte non vide $B_j$ avec $1 \leqslant j \leqslant 5$ ; ôter un jeton de la
  boîte $B_j$ et ajouter deux jetons à la boîte $B_{j+1}$.
  \item[\emph{Type} 2 :] Choisir une boîte non vide $B_k$ avec $1 \leqslant k \leqslant 4$ ; ôter un jeton de la
  boîte $B_k$ et échanger les contenus des boîtes (éventuellement vides) $B_{k+1}$ et $B_{k+2}$.
\end{itemize}
Est-il possible, à la suite d'un nombre fini de telles opérations, que les boîtes $B_1, B_2, B_3, B_4, B_5$
soient vides et que la boîte $B_6$ contienne exactement $2010^{2010^{2010}}$ jetons ? (Noter que
$a^{b^c} = a^{(b^c)}$.)
